\documentclass[12pt]{article}
\usepackage[utf8]{inputenc}
\usepackage[T1]{fontenc}
\usepackage{hyperref}
\title{Viti Language Resources for K-12 Curriculum Development}
\author{Compiled by Claude Sonnet 5 (claude-sonnet-5) for Andr\'as Kornai}
\date{}
\begin{document}
\maketitle

\section{Introduction}
Viti (ISO 639-3 vit; also Biti, V\"ot\"o or B\"ot\"o) is a language of Nigeria. A Wikipedia summary places it in Taraba State, in the village of Antere near the Cameroon border, and classifies it as Narrow Grassfields. The brief for this survey describes it as a stable Atlantic-Congo language of Nigeria. The two descriptions may not agree, so check the classification and location against Glottolog (viti1241) before relying on either. The search found no curriculum materials, primers, or school programmes for Viti, and no teaching materials of any kind. This dossier is therefore short.

\section{Primary Educational Resources}
None found in the search results. No primers, textbooks, or school programmes were found.

\section{Linguistic and Reference Materials}
\begin{itemize}
\item Wikipedia, Viti language: \url{https://en.wikipedia.org/wiki/Viti_language}
\item Endangered Languages Project, Viti bibliography (the page was not opened): \url{https://www.endangeredlanguages.com/node/38113/bibliography}
\item Glottolog reference entry, Roger Blench 2020 (not opened; its relevance to Viti is unconfirmed): \url{https://glottolog.org/resource/reference/id/586182}
\end{itemize}

\section{Religious and Community Resources}
\begin{itemize}
\item Joshua Project, Viti people group in Nigeria: \url{https://joshuaproject.net/people_groups/20770}
\item Joshua Project, Viti language page: \url{https://joshuaproject.net/languages/vit}
\end{itemize}
No scripture translation was confirmed in the search results.

\section{Digital and Audio Resources}
\begin{itemize}
\item Global Recordings Network, language page for vit (its contents were not verified): \url{https://globalrecordings.net/en/language/vit}
\end{itemize}

\section{Curriculum Development Context}
There is almost no published educational material for Viti. A K-12 programme would have to start with community-led documentation, an agreed orthography, and recordings, working with Viti speakers. The classification and location question above should be settled first.
\end{document}